\documentclass[11pt,a4paper]{article}

\usepackage[english]{babel}
\usepackage[T1]{fontenc}
\usepackage{lmodern}
\usepackage[a4paper,margin=1in]{geometry}

\usepackage{xcolor}
\usepackage{graphicx}
\usepackage{float}
\usepackage{booktabs}
\usepackage{array}
\usepackage{caption}
\usepackage{subcaption}
\usepackage{tabularx}

\usepackage{amsmath,amssymb,bm}
\usepackage{enumitem}
\usepackage{titlesec}
\usepackage[authoryear]{natbib}
\usepackage{hyperref}
\usepackage{fancyhdr}
\usepackage{pdflscape}
\usepackage{xfrac}

\usepackage{tikz}
\usetikzlibrary{arrows.meta, positioning, shapes.geometric}

\definecolor{darkblue}{rgb}{0,0,0.8}
\hypersetup{
	colorlinks=true,
	citecolor=darkblue,
	linkcolor=darkblue,
	urlcolor=darkblue
}

\setlength{\parskip}{0.55em}
\setlength{\parindent}{0pt}
\begin{document}
	

\tableofcontents

\pagebreak

\section{Introduction}

% ============================================================
% Essay 3 — Section B: Financial Subordination and World Money
% Draft v4 — 2026-03-23
% ============================================================
% Compile with: natbib (author-year), bibentry sectionB_references.bib
% This file is a standalone section meant to be \input{} into the main document.
% ============================================================

\subsection{Financial Subordination and World Money}
\label{sec:bridge}

\subsubsection{Historical transit: commodity money to credit money}
\label{sec:B1}

Money does not originate as credit. It originates as the commodity form taken by the universal equivalent---the one commodity whose use-value is to be the social expression of the value of all others \citep[\S3.1]{Basu2022}. In Marx's account, the exchange process itself produces a single commodity whose social function is to serve as the general measure of value and the universally accepted means of payment. Gold occupies this position not because of a natural property but because a historically specific process of social validation singles it out: the function is socially determined, the form is historically contingent \citep[\S3.2.1]{Basu2022}. What matters for theoretical purposes is not gold but the universal equivalent as a category---the socially validated device through which privately produced commodities are recognised as bearers of abstract labour.

Marx traces the developmental sequence by which credit money emerges from this commodity foundation, but the analysis remains at a level of abstraction he did not fully close. \citet[\S3.2.3]{Basu2022} reconstructs the chain. The bill of exchange, issued in the course of commercial transactions, is discounted by a bank; the bank issues its own liability against the discounted bill; the central bank stands behind commercial bank liabilities as the bank of banks; the state makes central bank liabilities legal tender. At each step the credit instrument substitutes for the commodity anchor, and at the final step social validity---the condition any form of money must satisfy to function as the universal equivalent---is guaranteed by the institutional authority of the state, not by the intrinsic value of a commodity \citep[\S3.2.3]{Basu2022}. The commodity anchor is no longer needed. What has changed is the form; what has not changed is the social function. Central bank liabilities and commercial bank deposits together now perform what gold once performed: the social materialisation of value, the mechanism through which privately produced commodities are validated as socially necessary.

The operative conversion factor linking value magnitudes to monetary magnitudes in a credit money system is the monetary expression of value---the ratio at which abstract labour-time is expressed in the monetary unit \citep{Foley1982, Basu2022}. What remains to be established is what this means when the universal equivalent crosses the border.


\subsubsection{World money under the monetary expression of value framework}
\label{sec:B2}

Marx's world money is bullion. When money leaves the sphere of national circulation it strips off its local garbs and returns to its original form---the commodity that functions as the absolute social materialisation of wealth \citep[p.~51]{Ivanova2013}. In Marx's formulation this is not a secondary function but the fullest expression of the potentialities and contradictions of the money-form: in the world market, money functions as the commodity whose natural form is the direct social form of realisation of human labour in the abstract. The question \citet[p.~45]{Ivanova2013} poses is whether this category can extend to a national fiat currency---whether a credit money instrument issued by a particular state can perform the function that Marx reserved for bullion.

Her answer turns on a distinction between substance and function. The fundamental feature of Marx's theory, she argues, is not its commodity-basis but its conception of money as the social expression of value, whose magnitude is determined by socially necessary labour embodied in the commodity. Money is an expression of the class relation between labour and capital, and the management of money is inseparable from the social form of imposition of work \citep[p.~45]{Ivanova2013}. If this is the theoretically load-bearing condition---not the physical substance of the universal equivalent but its social function---then the form of the universal equivalent is contingent. On this reading, the dollar is the closest equivalent to genuine world money the capitalist world market has known \citep[p.~52]{Ivanova2013}, not because it is a commodity but because it performs the function: it is the universally accepted means of international payment, reserve asset, and unit of account.

\citet[p.~6]{Ivanova2020} specifies the condition under which this claim holds. Credit money becomes money---income rather than mere promise---only if it enters the domain of social production as the embodiment of the value of social labour and social purchasing power. Central bank liabilities acquire the attributes of the universal equivalent not by state decree alone but through their social validation in production and circulation. The dollar's world money status is therefore not a legal fact but a social relation, continuously reproduced through the global circuit of capital. This condition is what separates Ivanova's reading from a merely nominalist account of fiat money.

The monetary expression of value framework provides a lens through which to assess this argument. The monetary expression of value, defined as the ratio of monetary value added to productive labour hours, is the operative conversion factor within any national monetary system \citep{Foley1982}. The Marxian exchange rate---the ratio of one country's monetary expression of value to another's---preserves social labour-time in exchange: it is the rate at which one hour of productive labour in economy~A exchanges for one hour in economy~B \citep[\S3.2.7]{Basu2022}. Ivanova's form-function argument is coherent with this framework. The dollar as world money means the US monetary expression of value functions as the implicit international reference against which peripheral monetary expressions are measured. This is not a formal derivation from the framework---it is a compatibility that allows taking her account at face value for its central claim.

What Ivanova's account does not fully theorise is the mechanism by which this position is structurally reproduced. Her analysis establishes that the dollar performs the world money function and that this performance is grounded in social relations of production, not in commodity substance or state fiat. But it does not explain why the dollar specifically retains this position, nor why peripheral subordination deepens rather than dissolves as US industrial dominance erodes. Ivanova identifies the explanandum; the explanans requires the institutional and structural account that follows.


\subsubsection{Bretton Woods and the peripheral balance-of-payments constraint}
\label{sec:B3}

The dollar-gold peg at \$35 per ounce, operative from 1944, institutionalised the dollar as the form of world money during the essay's period of interest. The United States occupied the apex of the monetary hierarchy not through market competition but by institutional design: it emerged from the war as the dominant creditor with substantial gold reserves, and the Bretton Woods architecture codified this position into the operating rules of the international monetary system \citep[p.~483]{Vasudevan2009}. The exorbitant privilege---the capacity to borrow internationally in one's own currency and to earn returns on foreign assets that exceed the cost of foreign liabilities---was operative from the outset, not only after the system's collapse \citep[p.~473]{Vasudevan2009}.

The asymmetry between centre and periphery under this architecture is structural, not contingent on any particular balance of trade. The country issuing world money can finance current account deficits in its own currency, and in doing so it centralises and recycles international surpluses through its financial system \citep[pp.~243--244]{Vasudevan2019}. It acts, in effect, as the banker to the world---borrowing short in the form of reserve holdings by other central banks and lending long in the form of foreign investment and military expenditure. Peripheral countries cannot do this. Their external liabilities are denominated in the dominant currency, and the settlement of international payments requires access to that currency. Instead of being able to shift the burden of adjustment onto trading partners, as the dominant country can, the periphery bears the brunt of debt-deflationary dynamics disproportionately \citep[p.~486]{Vasudevan2009}.

The periphery therefore operates under a balance-of-payments-constrained circuit of capital \citep{Vasudevan2016, Abeles2013}. Domestic credit creation is semi-endogenous---not in the sense that the central bank has no discretion over monetary policy, but in the sense that the ceiling on sustainable domestic monetary expansion is set by the accumulation of reserves denominated in the dominant currency, by export earnings, and by the terms on which external borrowing is available. The peripheral central bank holds the dominant country's liabilities as the necessary form of international liquidity, with no symmetric obligation on the other side. Reserve accumulation under these conditions is not a policy choice. It is the structural form through which the peripheral monetary system is anchored to the international monetary hierarchy---a form of financial subordination that reproduces the asymmetry regardless of domestic policy orientation \citep[p.~244]{Vasudevan2019}.

Chile's monetary system during the Bretton Woods period operates within this constraint. The closing of the gold window in 1971 dissolved the nominal anchor without dissolving the dollar's world money function \citep[p.~486]{Vasudevan2009}---a structural break whose implications for the transmission mechanism are tested through the post-1973 placebo specification in Part~C.


\subsubsection{Financial subordination as a state variable}
\label{sec:B4}

The ratio of exchange-rate-adjusted international reserves to the domestic monetary base---LevR---is the balance sheet expression of this position. It encodes not a domestic policy variable but the structural coverage of the peripheral central bank's monetary liabilities by the international universal equivalent---the ratio at which the national monetary system is anchored to the world monetary system during the Bretton Woods period. Section~C formalises this position as a state variable governing the transmission of external monetary shocks into domestic distributive dynamics.


\pagebreak
% REF References
\label{references}
\bibliographystyle{apalike}
\bibliography{references}


\end{document}
